\section{Introduction}
\label{sec:intro}

Parkinson's disease (PD) is the second most common neurodegenerative disorder,
and disturbed gait is one of its earliest and most visible motor signs. Clinical
scales such as the UPDRS rely on expert observation and are coarsely quantised,
which is why sensor-based measures are attractive. Instrumented insoles record
the vertical ground reaction force (VGRF) under each foot cheaply and at high
rate, and the resulting series carries the stride-timing irregularity,
shortened swing phase and left--right asymmetry that characterise parkinsonian
gait \cite{hausdorff2007, frenkel2005treadmill}.

Accuracies published on the PhysioNet gait database span roughly 70--99\,\%,
much wider than the differences one would expect between the model families
involved. A two-minute recording yields more than a hundred overlapping
windows, and if those are divided at random then windows of one subject end up
in both training and test. A classifier can then succeed by recognising the
individual instead of the disease -- a textbook case of leakage
\cite{kaufman2012leakage}, and a failure mode repeatedly documented in medical
machine learning \cite{varoquaux2022machine}. This report therefore asks two
questions at once: how well a state-space model classifies these signals, and
how much of a reported accuracy is an artefact of the split.

\subsection{Related work}
\label{ssec:sota}
Classical pipelines on this database pair hand-crafted statistics of the VGRF
channels with Random Forests or kernel SVMs. Deep models -- 1D-ResNets
\cite{he2016resnet}, recurrent networks and more recently transformers -- take
either the raw signal or sequences of feature vectors. Structured state-space
models are a newer alternative: S4 \cite{gu2022s4} and then Mamba
\cite{gu2023mamba} reach transformer-level sequence modelling at linear rather
than quadratic cost in sequence length, because Mamba makes the state transition
depend on the input instead of being fixed. The architecture has been applied to
speech, EEG and ECG; we are not aware of a VGRF-based PD study that uses it
under a subject-disjoint protocol.

\subsection{Contributions of this work}
\label{ssec:contr}
We train a Mamba classifier on this dataset and compare it with the Random
Forest, SVM-RBF, 1D-ResNet and LSTM baselines from the previous stage of the
project, evaluating every model with the same code and the same subject-disjoint
folds. Two findings came out of the experiments rather than out of the plan. The
input representation mattered more than the architecture: feeding raw samples
made training collapse on two of five folds, and describing each second by its
channel statistics both removed the collapse and improved accuracy. And the
split unit mattered more than either: re-running the identical pipeline with
window-wise splits raises Random Forest accuracy from 73\,\% to 98\,\%, which we
report as a caution rather than as a result.
